\documentclass[12pt]{article}
\usepackage[T1]{fontenc}
\usepackage[utf8]{inputenc}
\usepackage{lmodern}
\usepackage{textcomp}
\usepackage{enumitem}
\usepackage{geometry}
\geometry{margin=1in}
\begin{document}
\begin{center}
Answer Key - Computer Science - Graduate Level Assessment 18 \
\small (Answers are ordered exactly as in the question paper.)
\end{center}

\section*{Multiple Choice}
\begin{enumerate}[label=\arabic*.]
\item \textbf{Answer:} B
\\ \textit{Options:} A) A procedure for testing libraries or other program components for vulnerabilities; B) Whole-system testing for security flaws and bugs; C) A security-minded form of unit testing that applies early in the development process; D) All of the above
\\ \textit{Note:} Correct option: B - Whole-system testing for security flaws and bugs
\item \textbf{Answer:} B
\\ \textit{Options:} A) Message Condentiality; B) Message Integrity; C) Message Splashing; D) Message Sending
\\ \textit{Note:} Correct option: B - Message Integrity
\item \textbf{Answer:} A
\\ \textit{Options:} A) the matrix, if stored directly, is large and can be clumsy to manage; B) it is not capable of expressing complex protection requirements; C) deciding whether a process has access to a resource is undecidable; D) there is no way to express who has rights to change the access matrix itself
\\ \textit{Note:} Correct option: A - the matrix, if stored directly, is large and can be clumsy to manage
\item \textbf{Answer:} C
\\ \textit{Options:} A) crypter; B) virus; C) backdoor; D) key-logger
\\ \textit{Note:} Correct option: C - backdoor
\item \textbf{Answer:} B
\\ \textit{Options:} A) Installing and configuring a IDS that can read the IP header; B) Comparing the TTL values of the actual and spoofed addresses; C) Implementing a firewall to the network; D) Identify all TCP sessions that are initiated but does not complete successfully
\\ \textit{Note:} Correct option: B - Comparing the TTL values of the actual and spoofed addresses
\end{enumerate}

\section*{True / False}
\begin{enumerate}[label=\arabic*.]
\setcounter{enumi}{5}
\item \textbf{Answer:} False
\\ \textit{Options:} A) True; B) False
\\ \textit{Note:} LLM-generated false statement. Correct answer: 'a read outside bounds of a buffer'.
\item \textbf{Answer:} False
\\ \textit{Options:} A) True; B) False
\\ \textit{Note:} LLM-generated false statement. Correct answer: 'Session layer'.
\item \textbf{Answer:} True
\\ \textit{Options:} A) True; B) False
\\ \textit{Note:} LLM-generated true statement based on MCQ.
\item \textbf{Answer:} True
\\ \textit{Options:} A) True; B) False
\\ \textit{Note:} LLM-generated true statement based on MCQ.
\item \textbf{Answer:} True
\\ \textit{Options:} A) True; B) False
\\ \textit{Note:} LLM-generated true statement based on MCQ.
\end{enumerate}

\section*{Long Form Response}
\begin{enumerate}[label=\arabic*.]
\setcounter{enumi}{10}
\item \textbf{Answer:} def even\_num(x): | return True | return False
\\ \textit{Note:} Students should provide a detailed explanation referencing algorithm steps.
\item \textbf{Answer:} def count\_elim(num): | return count\_elim
\\ \textit{Note:} Students should provide a detailed explanation referencing algorithm steps.
\end{enumerate}

\vfill
\noindent\textit{End of Answer Key}
\end{document}